\section{Les Boucles}

\subsection{Boucle while (Tant que)}

Conçue pour itérer tant qu'une condition reste vraie. Le nombre de répétitions \textbf{n'est pas connu à l'avance}.

\begin{verbatim}
while [condition]:
    # Bloc exécuté tant que True
    [instructions]
    [mise_à_jour_de_la_condition] 
\end{verbatim}

Il est impératif de modifier une variable intervenant dans la condition à l'intérieur du bloc, sinon la boucle tournera à l'infini!

\subsection{Boucle for (Pour chaque)}

Conçue pour itérer sur une \textbf{séquence connue} (liste, chaîne de caractères, dictionnaire, etc.). La boucle s'arrête automatiquement à la fin de la séquence.

\begin{verbatim}
for [var_temporaire] in [séquence]:
# Bloc exécuté une fois par élément
    [instructions_avec_var_temporaire]
\end{verbatim}

À chaque tour, la var\_temporaire prend automatiquement la valeur de l'élément suivant dans la séquence.

\subsubsection{L'outil \texttt{range()}}

Génère une séquence de nombres entiers. Idéal pour répéter une action un nombre de fois précis.

\begin{itemize}
\item \texttt{range(fin)} : De \texttt{0} jusqu'à \texttt{fin - 1}.
\item \texttt{range(début, fin)} : De \texttt{début} jusqu'à \texttt{fin - 1}.
\item \texttt{range(début, fin, pas)} : De \texttt{début} jusqu'à \texttt{fin - 1}, en avançant selon le \texttt{pas}.
\end{itemize}

\subsubsection{Mots-clés de Contrôle}

Ces mots-clés s'utilisent à l'intérieur du bloc indenté des boucles \texttt{for} ou \texttt{while} (souvent combinés avec un \texttt{if}).

\begin{itemize}
\item \texttt{break} : Arrête l'itération actuelle et sort définitivement de la boucle entière.
\item \texttt{pass} : Ne fait rien (bien visuelment)
\end{itemize}